% Generated by report/tools/tables.py from the result files of the repository. Do not edit.
\begin{table}[tbp]
\centering
\footnotesize
\setlength{\tabcolsep}{3pt}
\caption[E5 : graines nécessaires pour retrouver le classement]{Expérience E5 : combien de graines retrouvent le bon classement ? Pour chaque comparaison de E2, 10\,000 tirages avec remise de $n$ essais par programme parmi ses essais (la distribution observée tient lieu de programme), puis la part des tirages dont les moyennes donnent le bon ordre, en pourcentage. Bon ordre : celui des moyennes de tous les essais choisis (10, ou 20 avec l'écart 20) ; à 10 graines, PPO PyTorch sur HalfCheetah tire dans les 30 essais de E1, comme préregistré. Différence : différence relative des deux moyennes, divisée par la plus grande en valeur absolue. \textbf{En gras}, les comparaisons serrées (différence sous 20~\%) ; la dernière ligne donne leur nombre et la médiane de leurs parts. Sources : \file{sim/analysis/e5.json}, \file{sim/analysis/e5\_s20.json}.}
\label{tab:e5}
\begin{tabular}{@{}llrrrrrrrrrr@{}}
\toprule
 & & \multicolumn{5}{c}{10 graines (résultat principal)} & \multicolumn{5}{c}{20 graines (écarts 12 et 20)} \\
\cmidrule(lr){3-7}\cmidrule(lr){8-12}
Tâche & Paire & écart & $n=3$ & 5 & 10 & 20 & écart & $n=3$ & 5 & 10 & 20 \\
\midrule
HalfCheetah & CrossQ / PPO PyTorch & 71 & 100 & 100 & 100 & 100 & 68 & 100 & 100 & 100 & 100 \\
 & CrossQ / PPO JAX & 72 & 100 & 100 & 100 & 100 & 67 & 98 & 100 & 100 & 100 \\
 & CrossQ / SAC & \textbf{8} & \textbf{56} & \textbf{60} & \textbf{64} & \textbf{70} & 26 & 73 & 79 & 89 & 96 \\
 & PPO PyTorch / PPO JAX & \textbf{3} & \textbf{73} & \textbf{71} & \textbf{77} & \textbf{84} & \textbf{4} & \textbf{47} & \textbf{51} & \textbf{56} & \textbf{58} \\
 & PPO PyTorch / SAC & 69 & 100 & 100 & 100 & 100 & 57 & 97 & 99 & 100 & 100 \\
 & PPO JAX / SAC & 70 & 100 & 100 & 100 & 100 & 56 & 96 & 99 & 100 & 100 \\
\midrule
Hopper & CrossQ / PPO PyTorch & 21 & 79 & 87 & 94 & 99 & \textbf{4} & \textbf{55} & \textbf{57} & \textbf{60} & \textbf{63} \\
 & CrossQ / PPO JAX & 27 & 87 & 94 & 99 & 100 & \textbf{8} & \textbf{61} & \textbf{65} & \textbf{71} & \textbf{80} \\
 & CrossQ / SAC & 32 & 83 & 89 & 96 & 99 & 38 & 91 & 96 & 99 & 100 \\
 & PPO PyTorch / PPO JAX & \textbf{7} & \textbf{65} & \textbf{69} & \textbf{76} & \textbf{85} & \textbf{4} & \textbf{58} & \textbf{59} & \textbf{63} & \textbf{68} \\
 & PPO PyTorch / SAC & 46 & 98 & 99 & 100 & 100 & 40 & 94 & 97 & 100 & 100 \\
 & PPO JAX / SAC & 50 & 99 & 100 & 100 & 100 & 43 & 95 & 99 & 100 & 100 \\
\midrule
Walker2d & CrossQ / PPO PyTorch & 57 & 100 & 100 & 100 & 100 & 52 & 99 & 100 & 100 & 100 \\
 & CrossQ / PPO JAX & 69 & 100 & 100 & 100 & 100 & 65 & 100 & 100 & 100 & 100 \\
 & CrossQ / SAC & 74 & 100 & 100 & 100 & 100 & 72 & 100 & 100 & 100 & 100 \\
 & PPO PyTorch / PPO JAX & 29 & 79 & 84 & 91 & 98 & 29 & 79 & 85 & 93 & 98 \\
 & PPO PyTorch / SAC & 39 & 89 & 95 & 99 & 100 & 41 & 93 & 97 & 100 & 100 \\
 & PPO JAX / SAC & \textbf{14} & \textbf{58} & \textbf{64} & \textbf{70} & \textbf{78} & \textbf{18} & \textbf{65} & \textbf{70} & \textbf{77} & \textbf{86} \\
\midrule
Swimmer & CrossQ / PPO PyTorch & \textbf{7} & \textbf{65} & \textbf{72} & \textbf{83} & \textbf{91} & \textbf{1} & \textbf{50} & \textbf{51} & \textbf{53} & \textbf{53} \\
 & CrossQ / PPO JAX & \textbf{4} & \textbf{74} & \textbf{82} & \textbf{91} & \textbf{97} & \textbf{15} & \textbf{84} & \textbf{91} & \textbf{98} & \textbf{100} \\
 & CrossQ / SAC & 57 & 100 & 100 & 100 & 100 & 50 & 99 & 100 & 100 & 100 \\
 & PPO PyTorch / PPO JAX & \textbf{11} & \textbf{81} & \textbf{88} & \textbf{96} & \textbf{100} & \textbf{14} & \textbf{89} & \textbf{94} & \textbf{99} & \textbf{100} \\
 & PPO PyTorch / SAC & 53 & 100 & 100 & 100 & 100 & 50 & 100 & 100 & 100 & 100 \\
 & PPO JAX / SAC & 58 & 100 & 100 & 100 & 100 & 57 & 100 & 100 & 100 & 100 \\
\midrule
\multicolumn{2}{@{}l}{Médiane, 24 comparaisons} &  & 94 & 97 & 99 & 100 &  & 93 & 97 & 100 & 100 \\
\multicolumn{2}{@{}l}{Médiane, comparaisons serrées} & 7 & 65 & 71 & 77 & 85 & 8 & 59 & 62 & 67 & 74 \\
\bottomrule
\end{tabular}
\end{table}
